\speciestitle{Hydrogen Chloride}{HCl}{HCl}{%
\item[Useful range:] 100\,--\,0.32\,hPa 
\item[Contact:] Lucien Froidevaux,
  \email{Lucien.Froidevaux@jpl.nasa.gov}}
\label{sec:standardHCl}

\subsection*{Introduction}

As in the previous MLS HCl data version, v2.2, the MLS \thisv retrievals
of the HCl standard product (from the 640\,GHz radiometer) use channels
from band 14, as a result of the deterioration observed since early 2006
in nearby band 13, originally targeted (with narrower channels than band
14) at the main HCl emission line center.  Full measurement days with
band 13 on from February 15, 2006, to the time of writing (May 2013) are
as follows: March~15, 2006 (2006d074), April~14, 2006 (2006d104),
January~6 through~8, 2009 (2009d006\,--\,2009d008), and January~24
through~27, 2010 (2010d024\,--\,2010d027). For days prior to February
16, 2006 and for the few days (as listed above) when band 13 is turned
on thereafter, the MLS Level~2 software also produces a separate
\texttt{HCl-640-B13} product (stored in the \texttt{L2GP-DGG} file),
using the band 13 radiances.  This product has slightly better precision
and vertical resolution in the upper stratosphere than the standard HCl
product.  The MLS team plans to turn band 13 on for a few days about
once every year or two (or maybe even less frequently) in order to
preserve its lifetime, estimated at a few days to a few weeks, based on
the channel counts and channel noise characteristics observed during the
3-day turn-on period in late January, 2010.  It is possible/likely that
this band will not be turned on again for a number of years.  Band 13
should provide better trend information for upper stratospheric data,
given its narrower channels.  Upper stratospheric trends from the
(uninterrupted 2004 to present) band 14 retrievals are too small,
compared to band 13 data and expectations (as well as versus ACE-FTS HCl
data).

\begin{figure}
  \centering\dontincludegraphics[angle=90,width=\linewidth]%
  {figures/HCl_2bands_0p46hpa_trends_2004_2010}
  \caption{Daily zonal averages for MLS HCl at 0.46\,hPa, from
    mid-August, 2004, through January, 2010, for the originally-targeted
    band 13 measurements (red points), now available only on occasion
    (to preserve lifetime), and the band 14 data (blue points).  The
    lines are simple linear fits through the daily data points; trend
    differences are apparent in this region of the atmosphere, where the
    information obtained from band 14 HCl data is not reliable enough.}
\label{fig:MLSHCl2bands}
\end{figure}


See Figure~\ref{fig:MLSHCl2bands} for an illustration of the trend
differences between these two MLS band measurements of upper
stratospheric HCl.  In the lower stratosphere, however, variations in
the two HCl products are closer together, and seasonal/geographical
variability are more pronounced.  We believe that the band 14 daily
global retrievals are completely suitable for use in studies of seasonal
and geographical variations (e.g., during polar winter/spring).

Table~\ref{tab:HClSummary} summarizes the MLS HCl resolution, precision,
and accuracy estimates as a function of pressure.  More discussion and
data screening recommendations for the MLS HCl \thisv data are provided
below.  Analyses describing detailed validation of the MLS (v2.2)
product and comparisons with other data sets are described in
\citet{FroidevauxEtal08HCl}.  Based on the fairly small overall changes
in \thisv HCl data (versus v2.2), the conclusions of the latter reference
should remain essentially unchanged.  Any minor updates will result from
new comparisons between MLS (\thisv) HCl and ACE-FTS HCl, which is also
being updated to a newer version (version 3).  We do not expect the
systematic uncertainty estimates in Table~\ref{tab:HClSummary} to change
significantly; however, an MLS team review of those estimates is
anticipated.

\subsection*{Changes from  v2.2}

While there were no large \thisv algorithmic changes relating to HCl, one
difference in the retrievals for HCl and other products derived from the
640-GHz MLS retrievals is that temperature information is now obtained
from the first retrieval phase (`Core'), as opposed to the 640-GHz
phases themselves; this led to overall improved efficiency, convergence,
and stability for the \thisv 640-GHz products.  A Level~1 change,
resulting from a small error in the spectral calibration files, which
led to all filter channel responses being shifted by a small fraction
(1$\%$) of the nominal channel widths, also had an impact on the HCl
results.  Mainly because of this change, \thisv MLS HCl abundances near
and just above the stratopause are a few percent less than the v2.2
abundances, and exhibit a steeper slope near very the top of the
recommended pressure range.  The slight oscillating behavior in HCl near
0.15 to 0.1\,hPa has led us to change the top boundary for recommended
HCl profiles to the MLS retrieval level at 0.32\,hPa.  This issue does
not seem to affect the band 13 MLS measurement of HCl, which can in
principle be used (on available measurement days) up to the 0.15\,hPa
level.

Other changes relating to the treatment of forward model radiance
continuum had an impact on species in the 640-GHz retrievals (mainly in
the lower stratosphere).  The background observed in the 640 GHz
radiances includes emissions from N\cs2, O\cs2, and H\cs2O.  There are
laboratory-based and ground-based models for the continuum absorptions
that are the basis for the MLS absorption model [{\it Pardo, 2001}, and
references therein].  These models were tested against MLS extinction
measurements from the wing channels in the 640 GHz radiometer; the
latitude dependence of this extinction was found to agree better with
the expected most plus dry continuum extinction values if the dry and
moist continuum functions were scaled by factors close to 20\%.  The
incorporation of this change improved the lower stratospheric retrievals
of most of the 640-GHz species (generally in terms of average negative
biases and their latitude dependence).

A comparison plot showing zonal average HCl contours and differences
between the two data versions for a typical month (April, 2006) is
provided in Figure~\ref{fig:MLSv2v3HClContour}.  For pressures larger
than or equal to 0.22\,hPa, the average differences between the two data
versions are typically within 0.1\,ppbv (a few percent).  The average
changes (globally within a few percent in most of this pressure range
for typical months) are within the estimated accuracy values (see
Table~\ref{tab:HClSummary}), which we have now changed (increased) to a
value of 10$\%$ (or about 0.3 ppbv) for pressures less than 10\,hPa,
given the trend issue for upper stratospheric HCl mentioned above.  The
largest percentage changes in HCl occur for very small mixing ratio
values; \thisv values can be larger than the v2.2 values by 20 to 50$\%$
(or more) under low HCl conditions in the lower stratosphere at low
latitudes or during winter at polar latitudes, even if these percentages
typically only reflect an increase of 0.1 ppbv (or less).  On occasion,
however, v2.2 zonal averages at 100\,hPa (mainly) or during southern 
hemisphere polar winter conditions at low altitude were slightly negative; 
this is no longer the case for \thisv data.  Whether the larger \thisv values 
at 100\,hPa (with averages now slightly above 0.1\,ppbv at low latitudes) 
are more realistic than v2.2 data remains to be seen, but this is a fairly 
minor change.  The precision estimated in the Level~2 files is essentially
unchanged from v2.2.

\begin{figure}
  \centering\dontincludegraphics[angle=90,width=\linewidth]{figures/HCl_v2_v3_Contour_2006m04}
\caption{Zonal averages for MLS HCl profiles during April, 2006, showing
  the MLS v2.2 HCl mixing ratio contours (top left panel), the \thisv
  contours (top right panel), and their differences in ppbv (\thisv minus
  v2.2, bottom left panel) and percent (\thisv minus v2.2 versus v2.2,
  bottom right panel).}
\label{fig:MLSv2v3HClContour}
\end{figure}

%\begin{figure}
%\centering\dontincludegraphics[angle=90,width=1.\linewidth]{figures/HCl...}
%\caption{Zonal average profiles ...}
%\label{fig:avk-HCl}
%\end{figure}

\subsection*{Resolution}
\onestandardavk{HCl}{HCl}%

Typical (rounded off) values for resolution are provided in
Table~\ref{tab:HClSummary}.  Based on the width of the averaging kernels
shown in Figure~\ref{fig:AvkHCl}, the vertical resolution for the
standard HCl stratospheric product is \texttildelow3\,km (2.7\,km at best in
the lower stratosphere), or about double the 640\,GHz radiometer
vertical field of view width at half-maximum; the vertical resolution
degrades to 4--6\, km in the lower mesosphere.  The along-track
resolution is \texttildelow200 to 350\,km for pressures of 2\,hPa or more, and
\texttildelow500\,km in the lower mesosphere.  The cross-track resolution is
set by the 3\,km width of the MLS 640\,GHz field of view.  The
longitudinal separation of MLS measurements, set by the Aura orbit, is
10\degsym\,--\,20\degsym\ over middle and lower latitudes, with much
finer sampling in polar regions.

\subsection*{Precision}

The estimated single-profile precision reported by the Level~2 software
varies from \texttildelow0.2 to 0.6\,ppbv in the stratosphere (see
Table~\ref{tab:HClSummary}), with poorer precision obtained in the lower
mesosphere.  These precision values have not changed significantly for
\thisv data.  The Level~2 precision values are often only slightly lower
than the observed scatter in the data, as evaluated from a narrow
latitude band centered around the equator where atmospheric variability
is often smaller than elsewhere, or as obtained from a comparison
between ascending and descending coincident MLS profiles.  The scatter
in MLS data and in simulated MLS retrievals (using noise-free radiances)
becomes smaller than the theoretical precision (given in the Level~2
files) in the upper stratosphere and mesosphere, where there is a larger
impact of a priori and smoothing constraints.  The HCl precision values
increase rapidly at pressures less than 0.2\,hPa, and are generally
flagged negative at pressures less than 0.1\,hPa; this indicates an
increasing influence from the \emph{a priori} (with poorer measurement
sensitivity and reliability).

\subsection*{Accuracy}

The accuracy estimates in the Table for v2.2 data came from a
quantification of the combined effects of possible systematic errors in
MLS calibration, spectroscopy, etc.\ on the HCl retrievals.  These
values are intended to represent 2 sigma estimates of accuracy.  For
more details, see the MLS validation paper by
\citet{FroidevauxEtal08HCl}.  For \thisv, however, given the trend issues
affecting the (band 14) standard HCl product in the upper stratosphere
and lower mesosphere, we now need to recommend a more conservative
accuracy estimate of 10\% in this region (or about 0.3 ppbv), rather
than the smaller numbers from the original (formal) estimates, which
should still apply to the (now very occasional) band 13 retrievals.
Given the better agreement between the two bands' retrievals in the
lower stratosphere, we maintain the formal accuracy estimates in this
region (see Table~\ref{tab:HClSummary}). Data users should be able to
reliably study seasonal and geographical changes in lower stratospheric
HCl (e.g., at high latitudes in winter or spring) with the current (band
14) standard HCl product.

\subsection*{Data screening}
\begin{description}
\item[Pressure range:] \screeningrule{100\,--\,0.32\,hPa}
  \standardrangestatement\ We note that the MLS values at 147\,hPa are
  are biased high, at least at low to mid-latitudes, and slightly more
  in the \thisv data than in the v2.2 data -- and these values are not
  recommended (particularly at low latitudes). Also, although the
  vertical range at the top end is recommended up to 0.32\,hPa, users
  should note the significant issues relating to HCl trend estimates in
  the upper stratosphere and lower mesosphere; average profiles in this
  region can be used for studies not involving trends (or accuracy
  requirements not as tight as 10\%).  The use of the band 13
  (intermittent) HCl data can/will continue to be carefully evaluated
  for trend-related issues.  
\standardprecisionrule \evenstatusrule

\item[Clouds:] \cloudsokrule

\item[Quality:] \qualityrule{1.2} This criterion removes profiles
  with the poorest radiance fits, typically significantly less than
  1$\%$ of the daily profiles. Results in this respect have improved, in
  comparison to v2.2 data.  For HCl (and for other 640\,GHz MLS
  products), this screening correlates well with the poorly converged
  sets of profiles (see below); we recommend the use of both the
  \texttt{Quality} and \texttt{Convergence} fields for data screening.
  The use of this screening criterion sometimes (but rarely) removes up
  to a few percent of global daily data (for example, during the first
  half of September, 2006, when some high latitude convergence and
  quality issues arose).

\item[Convergence:] \convergencerule{1.05} For the vast majority
  of profiles (99\% or more for most days), this field is less than
  1.05.  Results in this respect have improved, in comparison to v2.2
  data.  Nevertheless, on occasion, sets of profiles (typically one or
  more groups of ten profiles, retrieved as a `chunk') have this
  \texttt{Convergence} field set to larger values.  These profiles are
  usually almost noise-free and close to the \emph{a priori} profile,
  and need to be discarded as non-converged.  The \texttt{Quality} field
  (see above) most often yields poorer quality values for these
  non-converged profiles. The use of this screening criterion sometimes
  (but rarely) removes up to a few percent of global daily data (for
  example, during the first half of September, 2006, when some high
  latitude convergence and quality issues arose).
\end{description}

\subsection*{Review of comparisons with other datasets}

\citet{FroidevauxEtal08HCl} provided results of generally good comparisons
between MLS HCl and other satellite, balloon, and aircraft measurements.
Both MLS and ACE-FTS HCl values are generally larger (by about 10 to 15\%) 
than the HCl values from HALOE, especially at upper stratospheric altitudes; 
this feature has not changed, overall, with the new data version(s) from 
both MLS and ACE-FTS.  MLS HCl at 147\,hPa is biased high versus WB-57 aircraft 
in-situ (CIMS) measurements (low to mid-latitudes); while this is still true 
for \thisv data, MLS data on this pressure level may be useful and accurate enough 
at high latitudes. 

\subsection*{Artifacts}
\begin{itemize} 
\item We do not recommend the use of the MLS HCl standard product (from
  band 14) in the upper stratosphere and lower mesosphere, in terms of
  detailed trend studies, for reasons mentioned above.  The MLS HCl
  global results from band 13, although very infrequent (after early
  2006), are observed (and expected) to be more reliable in this
  respect.
\item The HCl values at 147\,hPa are biased high and generally not
  usable (except possibly at high latitudes).  Please consult the MLS
  team for further information.
\item Users should screen out the non-converged and poorest quality HCl
  profiles, as such profiles (typically a very small number per day)
  tend to behave unlike the majority of the other MLS retrievals.  See
  the criteria listed above.
\end{itemize}  

\begin{table}
\caption{Summary for MLS hydrogen chloride}
\vspace{\tablecaptionsep}
\label{tab:HClSummary}
\begin{minipage}{\linewidth}
\centering\begin{tabular}{ccccccm{25ex}}
\toprule
\parbox[m]{15mm}{\centering\textbf{Pressure}} &
\multicolumn{2}{c}{\parbox[m]{15mm}{\centering\textbf{Precision
      \footnote{Precision (1 sigma) for individual profiles; note that
        $\%$ values tend to vary strongly with latitude in the
        lower stratosphere.}}}} &
\parbox[m]{24mm}{\centering\textbf{Resolution\\ V $\times$ H}} &
\multicolumn{2}{c}{\parbox[m]{15mm}{\centering\textbf{Accuracy%
      \footnote{2 sigma estimate from systematic uncertainty
        characterization tests (but see text for estimates at 
        pressures lower than 10\,hPa); note that percent values tend 
        to vary strongly with latitude and season in the lower 
        stratosphere, due to the variability in HCl.}\\}}} 
&
\textbf{Comments} \\
\textbf{hPa} & \textbf{ppbv} & \textbf{\%} & \textbf{km} & \textbf{ppbv} & \textbf{\%} & \\
\midrule
0.2            & --- & ---             & ---              & ---   & ---     & Unsuitable for scientific use \\  
0.5            & 0.7 &   20            &   5 $\times$ 400 & 0.3   &   10    & Unsuitable for trend studies \\
1              & 0.5 &   15            &   4 $\times$ 300 & 0.3   &   10     & Unsuitable for trend studies \\
2              & 0.4 &   15            &   3 $\times$ 250 & 0.3   &   10     & Unsuitable for trend studies \\
5              & 0.3 &   10            &   3 $\times$ 200 & 0.3   &  10     & Unsuitable for trend studies \\
10             & 0.2 &   10            &   3 $\times$ 200 & 0.2   &  10     & \\
20             & 0.2 &   15            &   3 $\times$ 200 & 0.1   &  10     & \\
46             & 0.2 & 10 to $>$ 40    &   3 $\times$ 250 & 0.2   &  10 to $>$ 40  & \\
68             & 0.2 & 15 to $>$ 80    &   3 $\times$ 300 & 0.2   &  10 to $>$ 80   &  \\
100            & 0.3 & 30 to $>$ 100   &   3 $\times$ 350 & 0.15  &  10 to $>$ 100  &  \\
147            & 0.4 & 50 to $>$ 100   &   3 $\times$ 400 & 0.3  &   50
to $>$ 100  & High bias at low lats. (use with caution elsewhere) \\  
\bottomrule
\end{tabular}
\end{minipage} 
\end{table}

%%% Local Variables: 
%%% mode: latex
%%% TeX-master: "v4-x-report"
%%% End: 
